\begin{abstract}
Legal case retrieval must handle long documents, sparse citation cues, and highly domain-specific language. A key motivation of this study is a simple pattern: when both models are limited to 512 tokens, SAILER, the legal encoder proposed by THUIR@COLIEE 2023, can outperform BM25, but once BM25 is allowed to use the full document it becomes stronger again than truncated SAILER. This suggests that long-text truncation may be one of the main reasons embedding-based retrieval models (dense retrieval models) still struggle to consistently surpass lexical retrieval in legal case retrieval. To test this hypothesis, we adopt ModernBERT-base as a long-context encoder, first apply domain-adaptive pretraining (DAPT), i.e., continued pretraining with masked language modeling (MLM), on a U.S. case-law corpus, and then perform contrastive supervised fine-tuning (SFT) on Task~1 of the Competition on Legal Information Extraction and Entailment (COLIEE) 2025 and 2026.

Our experiments are constrained by a single NVIDIA RTX 4080 SUPER 16GB GPU, so the formal system uses 4096 tokens. In addition, under this hardware budget, the DAPT stage runs for about 60 hours and covers only 0.384 epoch, which means the model does not see the full 6.73 million-document U.S. case-law corpus. Even under this setting, three findings are clear. First, full-document queries are much better than citation-context queries: on 2026 validation, ModernBERT DAPT+SFT rises from 0.1038 to 0.1650 F1@5, and SAILER also rises from 0.1022 to 0.1213, when the query view switches from citation-context extraction to the full cleaned judgment. Second, static BM25 hard negatives make training nearly fail; stable convergence depends on dynamic negative sampling, while query-specific year scope filtering provides a smaller additional gain. Third, after adding scope filtering, LightGBM learning to rank, and post-processing whose key component is dynamic cut-off, the final FLNLP submission reaches 0.2935 F1 on COLIEE 2026 Task~1 and ranks 7th among 22 teams. Overall, this study shows that long-context handling, query representation, negative selection, and downstream ranking integration are the practical factors that determine whether legal embedding-based retrieval succeeds. The implementation of our method can be found at \url{https://github.com/emmanuelwithme/FLNLP-COLIEE2026}.
\end{abstract}

\begin{CCSXML}
<ccs2012>
   <concept>
       <concept_id>10002951.10003317.10003338</concept_id>
       <concept_desc>Information systems~Retrieval models and ranking</concept_desc>
       <concept_significance>500</concept_significance>
       </concept>
   <concept>
       <concept_id>10002951.10003317.10003338.10003343</concept_id>
       <concept_desc>Information systems~Learning to rank</concept_desc>
       <concept_significance>500</concept_significance>
       </concept>
   <concept>
       <concept_id>10010147.10010178.10010179</concept_id>
       <concept_desc>Computing methodologies~Natural language processing</concept_desc>
       <concept_significance>500</concept_significance>
       </concept>
   <concept>
       <concept_id>10002951.10003317.10003338.10003341</concept_id>
       <concept_desc>Information systems~Language models</concept_desc>
       <concept_significance>500</concept_significance>
       </concept>
 </ccs2012>
\end{CCSXML}

{\begin{sloppypar}
\ccsdesc[500]{Information systems~Retrieval models and ranking}
\ccsdesc[500]{Information systems~Learning to rank}
\ccsdesc[500]{Computing methodologies~Natural language processing}
\ccsdesc[500]{Information systems~Language models}
\end{sloppypar}}

\keywords{legal case retrieval, noticed case retrieval, embedding-based retrieval, dense retrieval, contrastive learning, ModernBERT, learning to rank}
